
\subsubsection{Limitations}

\textbf{L1: Single Implementation Approach} - We tested Hutchinson with 10 probes and power iteration with 5 iterations. Alternative configurations (50-100 probes, 20-50 iterations) or different estimators (Lanczos, randomized SVD) may exhibit different behavior.

\textbf{L2: Scale} - Experiments used GPT-2 small (125M) for 5000 steps (~320M tokens). Larger models or longer training may behave differently, though the severity of failure (zero measurements, 77,065\% perplexity deviation) suggests fundamental issues rather than scale-dependent phenomena.

\textbf{L3: Adaptive Lambda} - Lambda decay from 0.01 to 0.0063 did not prevent loss imbalance. More sophisticated tuning approaches may help, though the negative losses suggest implementation bugs rather than tuning issues.

\subsubsection{Alternative Approaches}

\textbf{Post-Hoc SVD Observational Study} - Measure stable rank via exact SVD on saved activations from pretrained checkpoints. Test correlations with LoRA rank, KV effective rank, and attention entropy. This approach tests observational correlation without requiring gradient-based control.

\textbf{Architectural Constraints} - Enforce low-rank structure via bottleneck layers, factorized attention, or explicit rank parameterizations. This avoids stochastic spectral estimation entirely.

\textbf{Gradient-Free Methods} - Decouple measurement from control by computing stable rank periodically via SVD without gradients, then applying explicit projection steps. This avoids backpropagating through spectral estimators.

\subsubsection{Implications}

This work documents that gradient-based Jacobian spectral control via Hutchinson trace and power iteration fails at the tested implementation level. The failure modes are specific: zero measurements from insufficient probes or gradient detachment, negative losses from autodiff pathologies, and training divergence from loss scale imbalance.

The gap between mathematical soundness and numerical feasibility is substantial. While J^T J ≈ F and Σ_{ℓ+1} ≈ J_ℓ Σ_ℓ J_ℓ^T are well-defined mathematically, operationalizing them via differentiable stochastic estimation in training loops exposes numerical instabilities.
